% Options for packages loaded elsewhere
\PassOptionsToPackage{unicode}{hyperref}
\PassOptionsToPackage{hyphens}{url}
\PassOptionsToPackage{dvipsnames,svgnames,x11names}{xcolor}
%
\documentclass[
  11pt,
]{article}
\usepackage{amsmath,amssymb}
\usepackage{iftex}
\ifPDFTeX
  \usepackage[T1]{fontenc}
  \usepackage[utf8]{inputenc}
  \usepackage{textcomp} % provide euro and other symbols
\else % if luatex or xetex
  \usepackage{unicode-math} % this also loads fontspec
  \defaultfontfeatures{Scale=MatchLowercase}
  \defaultfontfeatures[\rmfamily]{Ligatures=TeX,Scale=1}
\fi
\usepackage{lmodern}
\ifPDFTeX\else
  % xetex/luatex font selection
\fi
% Use upquote if available, for straight quotes in verbatim environments
\IfFileExists{upquote.sty}{\usepackage{upquote}}{}
\IfFileExists{microtype.sty}{% use microtype if available
  \usepackage[]{microtype}
  \UseMicrotypeSet[protrusion]{basicmath} % disable protrusion for tt fonts
}{}
\makeatletter
\@ifundefined{KOMAClassName}{% if non-KOMA class
  \IfFileExists{parskip.sty}{%
    \usepackage{parskip}
  }{% else
    \setlength{\parindent}{0pt}
    \setlength{\parskip}{6pt plus 2pt minus 1pt}}
}{% if KOMA class
  \KOMAoptions{parskip=half}}
\makeatother
\usepackage{xcolor}
\usepackage[margin=1in]{geometry}
\usepackage{longtable,booktabs,array}
\usepackage{calc} % for calculating minipage widths
% Correct order of tables after \paragraph or \subparagraph
\usepackage{etoolbox}
\makeatletter
\patchcmd\longtable{\par}{\if@noskipsec\mbox{}\fi\par}{}{}
\makeatother
% Allow footnotes in longtable head/foot
\IfFileExists{footnotehyper.sty}{\usepackage{footnotehyper}}{\usepackage{footnote}}
\makesavenoteenv{longtable}
\setlength{\emergencystretch}{3em} % prevent overfull lines
\providecommand{\tightlist}{%
  \setlength{\itemsep}{0pt}\setlength{\parskip}{0pt}}
\setcounter{secnumdepth}{-\maxdimen} % remove section numbering
\ifLuaTeX
  \usepackage{selnolig}  % disable illegal ligatures
\fi
\IfFileExists{bookmark.sty}{\usepackage{bookmark}}{\usepackage{hyperref}}
\IfFileExists{xurl.sty}{\usepackage{xurl}}{} % add URL line breaks if available
\urlstyle{same}
\hypersetup{
  pdftitle={Dual-theory reformulation of hydrogen-like precision spectroscopy},
  pdfauthor={Trey Morris with Claude Opus 4.7},
  pdfsubject={Introductory note for a new experimental collaborator in
optical/microwave atomic spectroscopy --- unified method for
hydrogen-like atoms + per-observable predictions and
precision-measurement comparisons},
  colorlinks=true,
  linkcolor={blue},
  filecolor={Maroon},
  citecolor={Blue},
  urlcolor={blue},
  pdfcreator={LaTeX via pandoc}}

\title{Dual-theory reformulation of hydrogen-like precision
spectroscopy}
\author{Trey Morris with Claude Opus 4.7}
\date{2026-05-28}

\begin{document}
\maketitle

\hypertarget{dual-theory-reformulation-of-hydrogen-like-precision-spectroscopy-intro-note-for-a-new-collaborator}{%
\section{Dual-theory reformulation of hydrogen-like precision
spectroscopy --- intro note for a new
collaborator}\label{dual-theory-reformulation-of-hydrogen-like-precision-spectroscopy-intro-note-for-a-new-collaborator}}

\textbf{Date:} 2026-05-28. \textbf{From:} Trey Morris (with Claude Opus
4.7). \textbf{Re:} Draft introductory email to a new experimental
collaborator in optical/microwave atomic spectroscopy. Tracks against
\href{https://github.com/temoTxt/PyPhysics/issues/88}{issue \#88}.
Methods-focused framing; honest-scope paragraph is load-bearing and must
be retained before sending. Source material: the Bethe--Salpeter
precision campaign
(\href{https://github.com/temoTxt/PyPhysics/issues/50}{\#50}), the
Li\(^{2+}\) Z-extension
(\href{https://github.com/temoTxt/PyPhysics/issues/78}{\#78}), and the
hydrogenic-ion g-factor Z-scan
(\href{https://github.com/temoTxt/PyPhysics/issues/82}{\#82}).

\begin{center}\rule{0.5\linewidth}{0.5pt}\end{center}

\hypertarget{draft-email}{%
\subsection{Draft email}\label{draft-email}}

\emph{Replace \texttt{{[}Name{]}}, the salutation, and the sign-off
before sending. Verify the numbers against the provenance table below.}

\begin{quote}
\textbf{Subject:} Dual-theory reformulation of hydrogen-like precision
spectroscopy --- would value your read

Dear {[}Name{]},

I'm a co-author on a reformulation of relativistic quantum mechanics ---
Gill's ``dual theory,'' built on a proper-time formulation --- and we've
been working through its predictions for precision atomic spectroscopy.
Since your work is in optical/microwave spectroscopy of hydrogen-like
systems, I'd value your read on the apparatus and your sense of which
measurements are most diagnostic. This note summarizes the method and
where it stands against measurement.

\textbf{The method, in one paragraph.} The dual theory recasts
relativistic dynamics through a proper-time canonical Hamiltonian,
\(K = H^2/(2mc^2) + mc^2/2\), with a dual-Dirac equation in place of the
standard Dirac equation. For precision atomic observables this reduces
to a clean structure: each \(g_s\)-dependent observable is the textbook
(Sommerfeld--Dirac / Fermi-contact) result times \((g_s/{-2})^{n}\),
with \(n=1\) for spin--orbit and Fermi-contact contributions and \(n=2\)
for two-fermion spin--spin couplings. The electron spin g-factor \(g_s\)
itself enters through a single cutoff parameter in the framework's
§III.D structure, and the whole apparatus scales across the
hydrogen-like sequence with the expected \(Z\)-powers. The upshot is
that one consistent prescription handles the Lamb shift, fine structure,
and hyperfine splitting across H, He\(^{+}\), Li\(^{2+}\), and up the
isoelectronic sequence.

\textbf{Where it stands against measurement (optical/microwave regime).}

\begin{longtable}[]{@{}
  >{\raggedright\arraybackslash}p{(\columnwidth - 6\tabcolsep) * \real{0.2500}}
  >{\raggedright\arraybackslash}p{(\columnwidth - 6\tabcolsep) * \real{0.2500}}
  >{\raggedright\arraybackslash}p{(\columnwidth - 6\tabcolsep) * \real{0.2500}}
  >{\raggedright\arraybackslash}p{(\columnwidth - 6\tabcolsep) * \real{0.2500}}@{}}
\toprule\noalign{}
\begin{minipage}[b]{\linewidth}\raggedright
Observable
\end{minipage} & \begin{minipage}[b]{\linewidth}\raggedright
Measurement
\end{minipage} & \begin{minipage}[b]{\linewidth}\raggedright
Framework
\end{minipage} & \begin{minipage}[b]{\linewidth}\raggedright
Note
\end{minipage} \\
\midrule\noalign{}
\endhead
\bottomrule\noalign{}
\endlastfoot
H Lamb shift \(2S_{1/2}\)--\(2P_{1/2}\) & \(1057.845(9)\) MHz &
\(\sim 1016\) MHz & Bethe-1947 estimate; \(\sim 42\) MHz (\(\sim 4\%\))
residual at the estimate's precision floor \\
H fine structure \(2P_{3/2}\)--\(2P_{1/2}\) & \(10\,969.13(10)\) MHz &
\(10\,962\) MHz & leading Dirac x anomalous-\(g\) \\
H 1S hyperfine (21-cm) & \(1\,420.405\,751\,768(2)\) MHz & \(1\,420.04\)
MHz & Fermi contact at the framework cutoff; \(\sim 0.4\) MHz
residual \\
Li\(^{2+}\) Lamb shift \(2S_{1/2}\)--\(2P_{1/2}\) & \(62\,765(21)\) MHz
(Schiffer 1995) & \(\sim 61\)--\(63\) GHz & \(Z^4\)-scaled
Bethe-estimate \\
Li\(^{2+}\) fine structure \(2P_{3/2}\)--\(2P_{1/2}\) & --- (no direct
precision measurement located) & \(\approx 888\) GHz & \(Z^4 \times\)
H \\
Li\(^{2+}\) 1s hyperfine & --- (theory comparator: Pachucki 2023) &
\(\approx 29.8\) GHz & \(Z^3 \times\) nuclear factors, \(I=3/2\) \\
\end{longtable}

\textbf{Honest scope.} At the precision the apparatus currently delivers
--- the leading-\(g_s\) / Bethe-estimate floor --- the framework
reproduces the standard QED results rather than predicting deviations
from them; the agreement is by construction of the reformulation, not an
independent corroboration. The frontier where a genuinely distinct
prediction could appear is a full proper-time one-loop calculation,
which we have not yet completed. I mention this up front because I'd
rather be precise about what's reformulation and what would be new
physics.

\textbf{What would be most useful from you.} Two things: (1) an expert
read on whether the proper-time / dual-Dirac reduction is faithful to
how these observables are actually measured and extracted; and (2) your
sense of which optical/microwave measurements --- higher-\(Z\)
hydrogen-like fine structure or hyperfine, improved \(2S\)--\(2P\)
intervals, or something else --- would most sharply test a reformulation
of this kind, particularly any regime where the proper-time
radiation-reaction structure might leave a measurable signature.

Happy to send the full per-observable derivations (each reproduced
symbolically in Mathematica) if useful. Thank you for considering it.

Best regards, {[}Trey Morris{]}
\end{quote}

\begin{center}\rule{0.5\linewidth}{0.5pt}\end{center}

\hypertarget{backing-provenance-table-for-trey-to-verify-before-sending-not-part-of-the-email}{%
\subsection{Backing provenance table (for Trey to verify before sending
--- not part of the
email)}\label{backing-provenance-table-for-trey-to-verify-before-sending-not-part-of-the-email}}

\begin{longtable}[]{@{}
  >{\raggedright\arraybackslash}p{(\columnwidth - 8\tabcolsep) * \real{0.2000}}
  >{\raggedright\arraybackslash}p{(\columnwidth - 8\tabcolsep) * \real{0.2000}}
  >{\raggedright\arraybackslash}p{(\columnwidth - 8\tabcolsep) * \real{0.2000}}
  >{\raggedright\arraybackslash}p{(\columnwidth - 8\tabcolsep) * \real{0.2000}}
  >{\raggedright\arraybackslash}p{(\columnwidth - 8\tabcolsep) * \real{0.2000}}@{}}
\toprule\noalign{}
\begin{minipage}[b]{\linewidth}\raggedright
Observable
\end{minipage} & \begin{minipage}[b]{\linewidth}\raggedright
Framework value
\end{minipage} & \begin{minipage}[b]{\linewidth}\raggedright
Source in repo
\end{minipage} & \begin{minipage}[b]{\linewidth}\raggedright
Measurement
\end{minipage} & \begin{minipage}[b]{\linewidth}\raggedright
Measurement source
\end{minipage} \\
\midrule\noalign{}
\endhead
\bottomrule\noalign{}
\endlastfoot
H Lamb shift \(2S_{1/2}\)--\(2P_{1/2}\) & \(\sim 1016\) MHz &
\href{../Quantum_Mechanics/Bethe_Salpeter/10_CrossComparison.md}{\texttt{Bethe\_Salpeter/10\_CrossComparison.md}}
PR E row & \(1057.845(9)\) MHz & CODATA-2018 \\
H fine structure \(2P_{3/2}\)--\(2P_{1/2}\) & \(10\,962\) MHz &
\href{../Quantum_Mechanics/Bethe_Salpeter/03_FineStructure.md}{\texttt{03\_FineStructure.md}}
BS-§14.2 & \(10\,969.13(10)\) MHz & Hagley \& Pipkin 1994 \\
H 1S hyperfine (21-cm) & \(1\,420.04\) MHz &
\href{../Quantum_Mechanics/Bethe_Salpeter/10_CrossComparison.md}{\texttt{10\_CrossComparison.md}}
§1 residual table & \(1\,420.405\,751\,768(2)\) MHz & NIST / Karshenboim
review \\
Li\(^{2+}\) Lamb shift & \(\sim 61\)--\(63\) GHz &
\texttt{12\_Li2plus\_Spectroscopy.md} §12.2 (branch
\href{https://github.com/temoTxt/PyPhysics/pull/85}{\#85}) &
\(62\,765(21)\) MHz & Schiffer et al., \emph{PRL} \textbf{74}, 2188
(1995) \\
Li\(^{2+}\) fine structure & \(887{,}920\) MHz &
\texttt{12\_Li2plus\_Spectroscopy.md} §12.3 (branch
\href{https://github.com/temoTxt/PyPhysics/pull/85}{\#85}) & none
located & --- (theory: \(Z^4\)-Dirac) \\
Li\(^{2+}\) 1s hyperfine & \(\sim 29.8\) GHz &
\texttt{13\_Li2plus\_Hyperfine.md} (branch
\href{https://github.com/temoTxt/PyPhysics/pull/86}{\#86}) & none direct
& Pachucki et al.~2023 (theory from Li\(^{+}\)) \\
\end{longtable}

\textbf{Note on the Li\(^{2+}\) rows:} the Li\(^{2+}\) predictions live
on the not-yet-merged PR branches
\href{https://github.com/temoTxt/PyPhysics/pull/85}{\#85} (Lamb shift +
fine structure) and
\href{https://github.com/temoTxt/PyPhysics/pull/86}{\#86} (hyperfine).
Verify against the merged versions once those PRs land, or against the
branch docs in the interim. The Li\(^{2+}\) fine-structure and hyperfine
rows have \textbf{no direct precision measurement} --- the email marks
them ``---'' honestly; do not imply an experimental comparison exists
for those two.

\textbf{Note on the honest-scope paragraph:} it is consistent with the
campaign's own framing in
\href{../Quantum_Mechanics/Bethe_Salpeter/10_CrossComparison.md}{\texttt{10\_CrossComparison.md}}:
\emph{``zero of the 28 recorded results constitute independent
corroborations of the dual theory's content distinct from textbook
QM.''} The hydrogenic-ion g-factor Z-scan
(\href{https://github.com/temoTxt/PyPhysics/issues/82}{\#82})
independently confirmed this on the Z-axis (the framework's cutoff is
Z-trivial; it inherits QED's bound-state structure rather than
predicting deviations). Keeping the honest-scope paragraph is
non-negotiable for an external-scientist audience.

\end{document}
